%%%%%%%%%%%%%%%%%%%%%%%%%%%%%%%%%%%%%%%%%%%%%%%%%%%%%%%%%%%%%%%%%%%%%%
% How to use writeLaTeX: 
%
% You edit the source code here on the left, and the preview on the
% right shows you the result within a few seconds.
%
% Bookmark this page and share the URL with your co-authors. They can
% edit at the same time!
%
% You can upload figures, bibliographies, custom classes and
% styles using the files menu.
%
%%%%%%%%%%%%%%%%%%%%%%%%%%%%%%%%%%%%%%%%%%%%%%%%%%%%%%%%%%%%%%%%%%%%%%

\documentclass[12pt]{article}

\usepackage{sbc-template}

\usepackage{graphicx,url}

%\usepackage[brazil]{babel}   
\usepackage[utf8]{inputenc}  

     
\sloppy

\title{Exercício 3 Sistemas Distribuídos}

\author{Christian Aguiar Plentz\inst{1}, João Roberto Lemos Fidellis\inst{2} }


\address{Ciência da Computação -- Universidade do Vale do Rio dos Sinos (UNISINOS)\\
  Caixa Postal 275 -- 93022-750 -- São Leopoldo -- RS -- Brazil
  \email{\{plentzchristian,JRFIDELLIS\}@edu.unisinos.br}
}


\begin{document} 

\maketitle

\section{Descreva vantagens e desvantagens dos três mecanismos de resolução de nomes estudado}
	\subsection{Navegação Iterativa}
		O cliente pede a um servidor DNS se ele tem um nome a ser resolvido e se não tiver
		o servidor de DNS manda outro servidor de DNS para o cliente tentar, uma desvantagem
		é que o cliente tem que fazer varias requisições sendo mais lento do lado dele
		mas não tem tanto esforço do lado dos servidores.
	\subsection{Navegação Recursiva}
		O cliente pede a um servidor DNS se ele tem um nome a ser resolvido e se nao tiver
		o servidor pede para outro servidor que se nao tiver pede a outro servidor e a resposta
		e passada de volta entre esses servidores, a vantagem é que o cliente só faz uma requisição
		em vez de fazer varias para outros servidores, dando assim qualidade de serviço para o cliente
	\subsection{Navegação Não-Recursiva}
		O cliente pede a um servidor DNS se ele tem um nome e se nao tiver o servidor pede para
		outros servidores resolverem o nome, o bom é que mantém somente uma requisição que o cliente
		deve fazer mas devido a isso toda a carga deve ser posta em um único servidor que procura
		o nome pelos outros servidores.
\section{Explique um cenário onde o sistema DynDNS pode ser útil}
	Em sistemas de vigilancia caseiro, pode ser que sua provedora te de um IP publico novo
	entao o DynDNS serve para te dar um nome para resolver para seu IP publico novo automaticamente
	para que nao precise por o IP publico para acessar seu sistema de vigilancia autonomo
3)Na parte de sistema de arquivos, qual a importância do VFS
	Oferecer uma API para novos sistemas de arquivos poderem suportar para ter compatibilidade
	com as livrarias de IO de um Sistema Operacional
\section{Ao analisar a ACL, qual o problema em termos um arquivo com permissão 777}
	se tiver 7 na parte do outros qualquer um pode ler e executar o arquivo gerando problemas
	de segurança no sistema um atacante não precisa ter acesso root para acessar um arquivo
	com permissão 7 na parte de outros.
\section{voce ve o sistema NFS como um RPC(Remote Procedure Call)? Explique.}
	Sim, para ter acesso ao Sistema de Arquivos através da rede é necessário
	invocar funções VFS pela rede para leitura e escritura dos arquivos.


\end{document}

